\section{LOOKUP 函数}
\subsection{VLOOKUP 函数}

需要在表格或区域中按行查找内容时，请使用 VLOOKUP。

语法：

=VLOOKUP(要查找的内容、要查找的位置、包含要返回的值的范围内的列号、返回表示为 1/TRUE 或 0/FALSE 的近似或精确匹配项)。

\begin{itemize}
    \item 查阅值应该始终位于所在区域的第一列，这样 VLOOKUP 才能正常工作。 
    
    例如，如果查阅值位于单元格 C2 内，那么您的区域应该以 C 开头。
    \item 返回值的列号：如果指定 B2:D11 作为区域，那么应该将 B 算作第一列，C 作为第二列，以此类推。
    \item 使用VLOOKUP将多个表中的数据合并到一个工作表上
    
    只要其中一个表具有与所有其他表共有的字段，就可以使用 VLOOKUP 将多个表合并为一个表。
\end{itemize}

\subsection{HLOOKUP 函数}

查找数组的首行，并返回指定单元格的值。

语法

HLOOKUP(lookup\_value, table\_array, row\_index\_num, [range\_lookup])

\begin{itemize}
    \item Lookup\_value    必需。 要在表格的第一行中查找的值。 Lookup\_value 可以是数值、引用或文本字符串。
    \item Table\_array    必需。 在其中查找数据的信息表。 使用对区域或区域名称的引用。
    \item Row\_index\_num    必需。 行号table\_array匹配值将返回的行号。
    \item Range\_lookup    可选。 一个逻辑值，指定希望 HLOOKUP 查找精确匹配值还是近似匹配值。 如果为 TRUE 或省略，则返回近似匹配值。
\end{itemize}

\subsection{XLOOKUP 函数}

XLookup函数是微软在2019年8月份发布的一个最新的查找函数。

XLOOKUP 函数搜索区域或数组，然后返回与它找到的第一个匹配项对应的项。 如果不存在匹配项，则 XLOOKUP 可以返回最接近 () 匹配项。

语法

=XLOOKUP(lookup\_value, lookup\_array, return\_array, [if\_not\_found], [match\_mode], [search\_mode]) 

\begin{itemize}
    \item lookup\_value （必需）要搜索的
    \item lookup\_array （必需）要搜索的数组或区域
    \item return\_array （必需）要返回的数组或区域
    \item $[\mathrm{if\_not\_found}]$（可选）如果找不到有效的匹配项，请返回你提供的 [if\_not\_found] 文本。
    
    如果找不到有效的匹配项，并且 [if\_not\_found] 缺失， 则返回\#N/A 。
    \item $[\mathrm{match\_mode}]$（可选）指定匹配类型：
        \begin{itemize}
            \item 0 - 完全匹配。 如果未找到，则返回 \#N/A。 这是默认选项。
            \item -1 - 完全匹配。 如果没有找到，则返回下一个较小的项。
            \item 1 - 完全匹配。 如果没有找到，则返回下一个较大的项。
            \item 2 - 通配符匹配。
        \end{itemize} 
    \item $[\mathrm{search\_mode}]$（可选）指定要使用的搜索模式：
        \begin{itemize}
            \item 1 - 从第一项开始执行搜索。 这是默认选项。
            \item -1 - 从最后一项开始执行反向搜索。
            \item 2 - 执行依赖于 lookup\_array 按升序排序的二进制搜索。 如果未排序，将返回无效结果。
            \item 2 - 执行依赖于 lookup\_array 按降序排序的二进制搜索。 如果未排序，将返回无效结果。
        \end{itemize} 
\end{itemize}

\section{文本函数}

\subsection{LEN 函数}

这个函数的主要功能，就是计算单元格中的字符数。

使用 LEN 函数计算字符数的时候，空格也是会被算作一个字符。

\begin{lstlisting}
A2 = 苹果apple
LEN(A2) = 7
\end{lstlisting}

\subsection{LENB 函数}

LENB 函数的主要作用，是计算单元格中的字节数量。

LENB 函数的计算结果大于 LEN 函数的计算结果，单元格中的所有汉字都被视为两个字节。

\begin{lstlisting}
    A2 = 苹果apple
    LENB(A2) = 9
\end{lstlisting}
    
\subsection{LEFT、LEFTB函数}

LEFT函数的主要功能，就是从左边开始返回特定数量的字符。

LEFTB 基于所指定的字节数返回文本字符串中的第一个或前几个字符。

\begin{itemize}
    \item Left函数把全角（如汉字）和半角字符（如数字与字母）都计作一个字符
    \item LeftB函数把全角字符计作两个字节、半角字符计作一个字节
  \end{itemize}

  语法：

  \verb|LEFT(text, [num_chars])|
  
  \verb|LEFTB(text, [num_bytes])|

\begin{lstlisting}
    A2 = 苹果apple
    LEFT(A2,2) = 苹果
    LEFT(A2,LENB(A2)-LEN(A2)) = 苹果
\end{lstlisting}

结合以上3个函数，可以实现将中英文字符分离。

\subsection{TEXTJOIN 函数}

TEXTJOIN 函数将多个区域和/或字符串的文本组合起来，并包括你在要组合的各文本值之间指定的分隔符。 如果分隔符是空的文本字符串，则此函数将有效连接这些区域。

语法：
\verb|TEXTJOIN(delimiter, ignore_empty, text1, [text2], …)|

\begin{itemize}
  \item delimiter：分隔符(必需)，文本字符串，或者为空，或用双引号引起来的一个或多个字符，或对有效文本字符串的引用。 如果提供一个数字，则将被视为文本。
  \item \verb|ignore_empty| (必需)，如果为 TRUE，则忽略空白单元格。
  \item \verb|text1|(必需)，要联接的文本项。 文本字符串或字符串数组，如单元格区域中。
  \item \verb|[text2， ...]|(可选)，要联接的其他文本项。 文本项最多可以包含 252 个文本参数 text1。 每个参数可以是一个文本字符串或字符串数组，如单元格区域。
\end{itemize}

例如，\verb|=TEXTJOIN(" ",TRUE,"The","sun","will","come","up","tomorrow")| 将返回The sun will come up tomorrow

\subsection{TEXTSPLIT 函数}

使用列和行分隔符拆分文本字符串。

TEXTSPLIT 函数的工作方式与文本转列向导相同，但采用公式形式。 它允许跨列拆分或按行向下拆分。 它是 TEXTJOIN 函数的反函数。 

语法：

\verb|=TEXTSPLIT(text,col_delimiter,[row_delimiter],[ignore_empty], |

\verb|[match_mode], [pad_with])|

TEXTSPLIT 函数语法具有下列参数：
\begin{itemize}
  \item text       要拆分的文本。 必需。
  \item \verb|col_delimiter|       标记跨列溢出文本的点的文本。
  \item \verb|row_delimiter|       标记向下溢出文本行的点的文本。 可选。
  \item \verb|ignore_empty|       指定 TRUE 以忽略连续分隔符。 默认为 FALSE，将创建一个空单元格。 可选。
  \item \verb|match_mode|    指定 1 以执行不区分大小写的匹配。 默认为 0，这会执行区分大小写的匹配。 可选。
  \item \verb|pad_with|           用于填充结果的值。 默认值为 \verb|#N/A|。
\end{itemize}